% ===== 底本 第 2 章《拉格朗日环面纤维化》；片段 tr02b；源块 B086--B104 =====

% [源 B086]
\subsection{等价}

% [源 B087]
拉格朗日纤维化是加在给定辛流形 $(M,\omega)$ 上的额外结构，我们还想理解两个这样的纤维化何时等价。

% [源 B088]
\begin{definition*}
两个正则拉格朗日纤维化 $\pi\colon M\to B$ 与 $\pi'\colon M'\to B'$ 称为\emph{同构的}（isomorphic），若存在纤维化辛同构（fibered symplectomorphism）$\Phi\colon M\to M'$，它覆盖微分同胚 $\phi\colon B\to B'$，即 $\Phi$ 是辛同构，且使下面的交换图交换：

% [源 B089]
\begin{tikzcd}
  M \arrow[r, "\Phi"] \arrow[d, "\pi"'] & M' \arrow[d, "\pi'"] \\
  B \arrow[r, "\phi"'] & B'
\end{tikzcd}
\end{definition*}

% [Tutor Note: B088 的“$\Phi$ 覆盖 $\phi$”等价于 $\pi'\circ\Phi=\phi\circ\pi$，即 B089 的方块交换；“fibered symplectomorphism” 指 $\Phi$ 把每条纤维 $\pi^{-1}(b)$ 同构地送到 $\pi'^{-1}(\phi(b))$，并满足 $\Phi^{\ast}\omega'=\omega$。]

% [源 B090]
\begin{exercise*}
不难证明 $\phi$ 必是整仿射流形的同构。
\end{exercise*}

% [Tutor Note: B090 的补出：B089 的交换性使 $\phi$ 把 $B$ 的周期格送到 $B'$ 的周期格，故 $\phi$ 在标准周期格坐标下的雅可比属于 $GL(n,\mathbb{Z})$；再由 B080/B081 的判定即得 $\phi$ 是整仿射同构。]

% [源 B091]
我们现在准备证明一个关于正则拉格朗日纤维化的“重构定理”，它类似于定理 [1.1]。

% [源 B092]
\begin{theorem*}
设 $M\xrightarrow{\pi} B$、$M'\xrightarrow{\pi'} B$ 是两个正则拉格朗日纤维化，分别带有整体拉格朗日截面 $\sigma$ 与 $\sigma'$。则存在唯一的纤维化辛同构 $\Phi\colon M\to M'$，使得下图

% [源 B093]
\begin{tikzcd}
  M \arrow[rr, "\Phi", dotted] \arrow[dr, "\pi", bend left=18] & & M' \arrow[dl, "\pi'", bend right=18] \\
  & B \arrow[ul, "\sigma", bend left=18] \arrow[ur, "\sigma'", bend right=18] &
\end{tikzcd}

% [源 B094]
交换。
\end{theorem*}

% [Tutor Note: B093 依 BRIEF §5 用 tikz-cd 重画：$M,M'$ 在上、$B$ 在下，$\Phi$ 为点线，$\pi,\pi',\sigma,\sigma'$ 四条边均为弧线（在 $M$ 与 $B$ 之间 $\pi$ 与 $\sigma$ 成透镜，在 $M'$ 与 $B$ 之间 $\pi'$ 与 $\sigma'$ 成透镜）。核对语义：$\pi\circ\sigma=\mathrm{id}_B$、$\pi'\circ\sigma'=\mathrm{id}_B$，而定理断言 $\Phi\circ\sigma=\sigma'$。]

% [源 B095]
\begin{proof}
用 Liouville 坐标构造 $\Phi$。
\end{proof}

% [Tutor Note: B095 的构造细节被原文略去，可补出如下：在局部 Liouville 坐标中两个纤维化都写成 $B\times\mathbb{T}^n$ 上的作用-角坐标，$\Phi$ 在每块上取恒等（至多差一个整仿射过渡），再用整体截面 $\sigma,\sigma'$ 固定纤维方向的平移，即得唯一性。]

% [源 B096]
事实上，在这些假设下我们可以显式地理解空间 $M$：$B$ 上的整仿射结构在余切丛 $T^{\ast} B$ 中诱导一个局部自由的格层（sheaf of lattices）$\Lambda$，于是我们可以自然地构造辛流形 $T^{\ast} B/\Lambda$，其投影 $T^{\ast} B/\Lambda\to B$ 的纤维是拉格朗日环面。注意 $T^{\ast} B/\Lambda\to B$ 总有由零截面给出的整体拉格朗日截面，因此重构定理 [2.20] 适用，推出 $M$ 与 $T^{\ast} B/\Lambda$ 纤维化辛同胚。

% [Tutor Note: B096 的构造可这样核对：$T^{\ast} B$ 带标准辛形式，$\Lambda$ 是整仿射结构给出的格层，沿 $\Lambda$ 取商后每条纤维是 $T_b^{\ast}B/\Lambda_b\cong\mathbb{T}^n$；若 $B$ 可缩使 $\Lambda$ 平凡，则 $T^{\ast}B/\Lambda\cong B\times\mathbb{T}^n$，与 B074 的局部模型一致（这也说明 $M$ 的局部模型为何只有这一种）。]

% [源 B097]
\begin{example*}
当 $M\to B$ 没有整体拉格朗日截面时，重构定理可能失效；下面是 [Eva23] 中的一个例子：考虑带坐标 $(t,x,y,z)$ 与辛结构 $dt\wedge dx+dy\wedge dz$ 的辛向量空间 $\mathbb{R}^4$。$\mathbb{Z}^3$-作用 $(a,b,c)\cdot (t,x,y,z)=(t,x+a,y+b,z+c)$ 保持辛结构，于是对这个 $\mathbb{Z}^3$-作用取商，得到辛流形 $\mathbb{R}\times\mathbb{T}^3$。再考虑 $\mathbb{R}\times\mathbb{T}^3$ 上的 $\mathbb{Z}$-作用 $1\cdot (t,x,y,z)=(t+1,x,y,x+y+z)$，它也保持辛结构，对这个 $\mathbb{Z}$-作用取商，得到辛流形 $(K,\omega)$。这就是著名的\emph{Kodaira--Thurston 流形}：它是紧辛流形，却不容许任何相容的复结构（有相容复结构才能成为 Kähler 流形）。

考虑投影 $K\to\mathbb{T}^2$，$(t,x,y,z)\mapsto (t,y)$。这是一个（主）正则拉格朗日环面纤维化，而有整体截面的主丛总是平凡的，即 $K\cong\mathbb{T}^4$。然而我们可以算出 $b_1 (K)=3$，而 $b_1 (\mathbb{T}^4)=4$，所以 $K\to\mathbb{T}^2$ 不容许任何整体拉格朗日截面。

另一方面，考虑带辛结构 $dt\wedge dx+dy\wedge dz$ 的标准环面 $\mathbb{T}^4$ 与投影 $\mathbb{T}^4\to\mathbb{T}^2$，$(t,x,y,z)\mapsto (t,y)$；它在 $\mathbb{T}^2$ 上诱导的整仿射结构与 $K\to\mathbb{T}^2$ 所诱导的相同。这说明：若不存在整体拉格朗日截面，重构定理可能失效。
\end{example*}

% [Tutor Note: B097 源文以 br 断行分成三段（Kodaira--Thurston 的构造、$K\to\mathbb{T}^2$ 无截面、$\mathbb{T}^4\to\mathbb{T}^2$ 的对照），译文按此分为三段，仍在同一例环境内。]
% [Tutor Note: B097 的“主丛若有整体截面则平凡，即 $K\cong\mathbb{T}^4$”应读作条件句：若 $K\to\mathbb{T}^2$ 存在整体拉格朗日截面，则主丛平凡从而 $K\cong\mathbb{T}^4$；而 $b_1(K)=3\neq4=b_1(\mathbb{T}^4)$ 恰说明这种截面不存在。原文句法省略，照译并以批注澄清。]
% [Tutor Note: B097 括注 “so that it becomes a Kähler manifold” 按字面直译会自相矛盾（无相容复结构恰是“不是 Kähler”的理由）；此处译为“有相容复结构才能成为 Kähler 流形”。依据：Kähler 流形按定义需相容可积复结构。]

% [源 B098]
\subsection{通量}

% [源 B099]
最后，还有另一种更几何的方式来写出作用-角坐标，即通过计算某些周期积分。对给定的拉格朗日纤维化 $M\to B$，假定 $M$ 是恰当的（exact），即 $\omega =d\lambda$，其中 $\lambda$ 是某个 $1$-形式；我们得到一个局部自由层 $H_1 (F_b;\mathbb{Z})\to\mathcal{S}\to B$，它沿万有覆盖 $p\colon\tilde{B}\to B$ 拉回后变为平凡。设 $c_1,\dotsb, c_n$ 是拉回丛 $p^{\ast}\mathcal{S}$ 的平凡化截面。

% [Tutor Note: B099 之所以要假定 $\omega=d\lambda$（$M$ 恰当），是为了让周期积分 $\int_{c_i}\lambda$ 有定义；若 $\lambda$ 换成 $\lambda+\eta$（$\eta$ 闭），则 $I_i$ 只差常数 $\frac{1}{2\pi}\int_{c_i}\eta$，而常数平移正是 B080 中整仿射变换 $T(b)=b+c$ 的平移部分，故定义在整仿射意义下无歧义（详见 B101）。]

% [源 B100]
\begin{definition*}
我们定义\emph{通量映射}（flux map）$I\colon\tilde{B}\to\mathbb{R}^n$ 为
\[
  I(\tilde{b})=(I_1 (\tilde{b}),\dotsb, I_n (\tilde{b}))=\left(\frac{1}{2\pi}\int_{c_1 (\tilde{b})}\lambda,\dotsb, \frac{1}{2\pi}\int_{c_n (\tilde{b})}\lambda\right).
\]
\end{definition*}

% [源 B101]
这个定义不依赖于本原 $1$-形式 $\lambda$ 的选取：不同的选取只相差一个常数平移，而常数平移是整仿射变换；它也不依赖于基 $\{c_i\}$ 的选取，因为任何两个选取相差 $GL(n,\mathbb{Z})$ 中的一个元素，这同样被整仿射结构吸收。

% [Tutor Note: B101 的“常数平移是整仿射变换”可取 $A=I_n\in GL(n,\mathbb{Z})$ 核对；而基的更换给出 $c_i'=\sum_j B_{ij}c_j$（$B\in GL(n,\mathbb{Z})$），于是 $I'=BI$，正是整仿射变换的线性部分。]

% [源 B102]
\begin{lemma*}
设 $\tilde{U}\xrightarrow{\cong} U$ 分别是 $\tilde{B}$ 与 $B$ 的开子集，则复合映射 $I\circ\left(p\vert_{\tilde{U}}\right)^{-1}\colon U\to\mathbb{R}^n$ 给出 $U$ 中的作用坐标。
\end{lemma*}

% [源 B103]
\begin{proof}
只需在 $M=U\times\mathbb{T}^n$ 带标准作用-角辛结构的情形验证。此时可写 $\lambda =\sum b_i dt_i$，且 $I_i (b)=\displaystyle\frac{1}{2\pi}\int_{c_i} b_idt_i =b_i$，其中 $c_i$ 是标准基。
\end{proof}

% [Tutor Note: B103 的 $\frac{1}{2\pi}\int_{c_i}b_i\,dt_i=b_i$ 依赖 $t_i$ 的周期为 $2\pi$（B066 的 standard 约定）：沿标准基圈 $c_i$，$\int dt_i=2\pi$，除以 $2\pi$ 后恰得 $b_i$。若把周期改成 $1$，此式须改为 $b_i/(2\pi)$。]

% [源 B104]
借助通量映射 $I\colon\tilde{B}\to\mathbb{R}^n$，我们把 $\mathbb{R}^n$ 上的标准整仿射结构沿 $I$ 拉回，便自然地得到 $\tilde{B}$ 上的整仿射结构。这个整仿射结构随后下降到 $B$，并给出与正则拉格朗日纤维化所给出的相同的 $B$ 上整仿射结构。

% [Tutor Note: B104 的“下降”依据：$I$ 沿万有覆盖的 Deck 变换只差整仿射变换（$I\circ\gamma=BI$，$B\in GL(n,\mathbb{Z})$，来自 B101），故 $\tilde{B}$ 上的整仿射结构在 $\pi_1(B)$ 作用下不变，从而下降到 $B$；它与 B085 的整仿射结构一致，由 B102（$I$ 局部给出作用坐标）保证。]
